% ---
% title: 重新审视扩散模型中的噪声：理论与实验
% date: 2026/3/28
% tag: Diffusion
% excerpt: DDPM 默认用白高斯噪声，但这篇论文问的是：换成彩色高斯或非高斯噪声会怎样。几个定理给出 MSE 下界和误差概率下界，实验里 FID 从 104 降到 70。
% ---
\begin{quote}
论文：\emph{Reconsidering Noise for Denoising Diffusion Probabilistic Models}（NeurIPS 2025 Workshop: SPIGM）
\end{quote}

\section{背景：白噪声的"默认地位"}

绝大多数去噪扩散概率模型（DDPM）在前向扩散和反向去噪时都默认用白高斯噪声（White Gaussian Noise）。理由很实际：分析方便，实现简单。至于它是不是最优选择，论文从理论出发回答了三个问题：

\begin{enumerate}
  \item 彩色高斯噪声或非高斯噪声能否替代白高斯噪声？
  \item 替换后 MSE 性能是否会改善？
  \item 有哪些理论依据支撑这种替换？
\end{enumerate}

\hrule

\section{核心理论结果}

\subsection{定理 1：去噪输出匹配输入的必要条件}

对于输入 \(x_0\)、去噪输出 \(\hat{x}_0\)，定义误差概率：

\[
p(e) = \Pr(\|x_0 - \hat{x}_0\| \geq \epsilon)
\]

利用 Fano 不等式，可推导出其下界：

\[
p(e) \geq \frac{\log(1-\bar{\alpha}_T) - 2}{\log(2\pi e)^N |K_{x_0}|} + 1
\]

要使 \(\hat{x}_0\) 有可能等于 \(x_0\)，上式右端必须 \(\leq 0\)。

推论：若输入数据的协方差矩阵行列式满足 \(|K_{x_0}| < \frac{1}{(2\pi e)^N}\)，则 \(x_0\) 与 \(\hat{x}_0\) 必然不同，完美重建从理论上就不可能实现。

\hrule

\subsection{定理 2：彩色高斯噪声的 MSE 更小}

DDPM 训练的核心目标是最小化噪声预测网络的 MSE：

\[
\min_{\hat{\epsilon}} \mathbb{E} \|\epsilon_t - \hat{\epsilon}_t\|_2^2
\]

根据估计误差与微分熵定理，该 MSE 的下界为：

\[
\text{MSE} \geq \frac{1}{(2\pi e)^N} e^{2h(\epsilon)}
\]

关键是彩色高斯噪声的最优协方差矩阵行列式更小（Hadamard 不等式），熵更低，MSE 下界更紧。

\begin{quote}
结论：在相同平均功率下，彩色高斯噪声比白高斯噪声能实现更低的 MSE。
\end{quote}

推论（Corollary 2）：非高斯噪声的熵低于同协方差矩阵的高斯噪声，所以在 MSE 最小化上更有优势。

\hrule

\subsection{定理 4：彩色噪声的误差概率下界更小}

进一步分析前向过程各时间步的误差概率，可以证明使用彩色高斯噪声时，误差概率的下界为：

\[
p_t(e) \geq 1 - \frac{2\bar{\alpha}_t}{\log(2\pi e)^N \bar{\beta}_t |K_n|}
\]

由于彩色噪声的 \(|K_n|\) 更大（相较于白噪声的对角矩阵），上式右端更小，即误差概率下界更低，理论上更有可能实现精确去噪。

\hrule

\subsection{定理 5：白噪声训练的 DDPM 可泛化至任意零均值对称噪声}

这一条最实用。

用白高斯噪声训练好的 DDPM，推理时即使初始噪声 \(x_T\) 来自任意零均值对称分布（均匀、拉普拉斯、余弦、Logistic 等），KL 散度仍然单调不增，\(t \to 0\) 时趋于 0。

证明用的是概率流 ODE 的 KL 耗散恒等式：

\[
\frac{d}{dt} \text{KL}(\mu_t \| q_t) = -g(t)^2 \underbrace{\int \mu_t \|\nabla \log \mu_t - \nabla \log q_t\|^2 \, dx}_{\geq 0}
\]

右端始终非正，故 KL 散度单调递减，且推导过程不依赖初始分布的具体形状，仅需其具有密度函数与有限二阶矩。

\begin{quote}
这个性质只对零均值对称分布成立。非对称分布（如 Gamma、指数分布）的实验结果显示无法正常生成图像。
\end{quote}

\hrule

\section{实验验证}

\subsection{数据集与设置}

\begin{center}
\begin{tabular}{lll}
\toprule
数据集 & 分辨率 & 用途 \\
\midrule
HuggingFace Butterflies（1000张） & 3×32×32 & 验证不同训练噪声下的 MSE 与 FID \\
LSUN Church-256 & 3×256×256 & 验证白噪声训练模型对不同推理噪声的泛化性 \\
\bottomrule
\end{tabular}
\end{center}

模型采用标准 2D-UNet 架构，T = 1000。

\subsection{MSE 对比（蝴蝶数据集）}

训练过程中，彩色高斯噪声的 MSE 收敛曲线始终低于白高斯噪声，均匀噪声居中，与理论预测一致。

\subsection{FID 对比（蝴蝶数据集）}

\begin{center}
\begin{tabular}{ll}
\toprule
噪声类型 & FID ↓ \\
\midrule
白高斯噪声 & 104.41 \\
均匀噪声 & 78.49 \\
\textbf{彩色高斯噪声} & \textbf{69.89} \\
\bottomrule
\end{tabular}
\end{center}

彩色高斯噪声生成的图像质量最高，FID 降低约 33\%。

\subsection{泛化实验（LSUN Church-256）}

在白噪声预训练模型上，用六种非高斯噪声进行推理：

均匀分布、拉普拉斯分布、余弦分布、Logistic 分布（零均值对称）都能生成高质量图像。Gamma 分布、指数分布（非对称）生成不出有效图像。

\hrule

\section{结论一览}

\begin{center}
\begin{tabular}{ll}
\toprule
结论 & 依据 \\
\midrule
彩色高斯噪声优于白高斯噪声 & MSE 更小，误差概率下界更低，FID 更优 \\
非高斯噪声（同协方差）优于高斯噪声 & 熵更低，MSE 下界更紧 \\
白噪声训练模型可泛化至零均值对称噪声 & KL 散度单调耗散定理 \\
非对称噪声不适用 & 缺乏理论保障，实验验证失败 \\
\bottomrule
\end{tabular}
\end{center}

这篇论文从信息论角度给 DDPM 的噪声选择补了理论依据：白高斯噪声不是唯一合理的选择，彩色高斯噪声的 MSE 更小、误差概率下界更低、FID 也更好，非对称噪声则没有理论保障。
